% GENERATED FILE. Edit canonical lessons, then run npm run generate:books.
% canonical-source-hash: fnv1a64:9ea605cfefce66fd
% canonical-lessons: HI-C478-phephra, HI-C478-dimag, HI-C478-nas, HI-C478-pasli, HI-C478-gurda

\chapter{Inside the Body}
\label{ch:hi-inside-the-body}
\hlchaptermodality{478}

\begin{chapteropening}
\textbf{By the end of this chapter:} \emph{I can name some organs inside the body.}

Complete the last lesson of chapter 478: I can name some organs inside the body.
\end{chapteropening}

\section[\texorpdfstring{\hi{फेफड़ा} (phephṛā)}{फेफड़ा (phephṛā)}]{\hi{फेफड़ा} (phephṛā) — a lung}

\label{lesson:HI-C478-phephra}



\begin{quote}
Before the new one: say the Hindi for a leg, then the Hindi for a heel.
\end{quote}

\subsection*{You'll want to know: \hi{फेफड़ा}}
\textbf{\hi{फेफड़ा}} — \emph{phephṛā} — \textquotedblleft{}a lung\textquotedblright{}.

\begin{quote}
\hi{गंदी} \hi{हवा} \hi{से} \hi{फेफड़ों} \hi{को} \hi{नुक़सान} \hi{होता} \hi{है।}
\end{quote}

\begin{quote}
\emph{gandī havā se phephṛoṁ ko nuqsān hotā hai.}
\end{quote}

\begin{quote}
Dirty air harms the lungs.
\end{quote}

Masculine.

\begin{cousinweb}
\textbf{\hi{फेफड़ा}} (\emph{phephṛā}) is a tadbhava: CDIAL lists Hindi \emph{phephṛā}, \textquotedblleft{}lungs\textquotedblright{}, under Sanskrit \emph{phupphusa}, \textquotedblleft{}lungs\textquotedblright{} (CDIAL 9090).
\end{cousinweb}

\begin{grammarlens}[title={What you've built}]
One more word for describing what you lived and felt.
\end{grammarlens}

\subsection*{Your turn}
\begin{itemize}
\raggedright
  \item \emph{Say it:} \emph{phephṛā}
  \item \emph{Say it:} \emph{phephṛā}, once more
  \item \emph{Say it:} \emph{gandī havā se phephṛoṁ ko nuqsān hotā hai.} — Dirty air harms the lungs.
  \item \emph{Recall:} say the Hindi for a leg, then the Hindi for a heel, then say \emph{phephṛā} again
\end{itemize}

\begin{tcolorbox}[breakable,colback=teal!4,colframe=teal!35!black,title={Before you move on}]
What does \hi{फेफड़ा} mean? (\textquotedblleft{}A lung\textquotedblright{}.) Say it once more.
\end{tcolorbox}
\section[\texorpdfstring{\hi{दिमाग़} (dimāġ)}{दिमाग़ (dimāġ)}]{\hi{दिमाग़} (dimāġ) — the brain; the mind}

\label{lesson:HI-C478-dimag}



\begin{quote}
Before the new one: say the Hindi for a heel, then the Hindi for a lung.
\end{quote}

\subsection*{You'll want to know: \hi{दिमाग़}}
\textbf{\hi{दिमाग़}} — \emph{dimāġ} — \textquotedblleft{}the brain; the mind\textquotedblright{}.

\begin{quote}
\hi{पढ़ाई} \hi{से} \hi{दिमाग़} \hi{तेज़} \hi{होता} \hi{है।}
\end{quote}

\begin{quote}
\emph{paṛhāī se dimāġ tez hotā hai.}
\end{quote}

\begin{quote}
Studying sharpens the mind.
\end{quote}

Masculine.

\begin{cousinweb}
\textbf{\hi{दिमाग़}} (\emph{dimāġ}) is usually traced to Arabic \emph{dimāġ}, \textquotedblleft{}brain\textquotedblright{}, which reached Hindi through Persian.
\end{cousinweb}

\begin{grammarlens}[title={What you've built}]
One more word for describing what you lived and felt.
\end{grammarlens}

\subsection*{Your turn}
\begin{itemize}
\raggedright
  \item \emph{Say it:} \emph{dimāġ}
  \item \emph{Say it:} \emph{dimāġ}, once more
  \item \emph{Say it:} \emph{paṛhāī se dimāġ tez hotā hai.} — Studying sharpens the mind.
  \item \emph{Recall:} say the Hindi for a heel, then the Hindi for a lung, then say \emph{dimāġ} again
\end{itemize}

\begin{tcolorbox}[breakable,colback=teal!4,colframe=teal!35!black,title={Before you move on}]
What does \hi{दिमाग़} mean? (\textquotedblleft{}The brain; the mind\textquotedblright{}.) Say it once more.
\end{tcolorbox}
\section[\texorpdfstring{\hi{नस} (nas)}{नस (nas)}]{\hi{नस} (nas) — a vein; a nerve}

\label{lesson:HI-C478-nas}



\begin{quote}
Before the new one: say the Hindi for a lung, then the Hindi for the brain.
\end{quote}

\subsection*{You'll want to know: \hi{नस}}
\textbf{\hi{नस}} — \emph{nas} — \textquotedblleft{}a vein; a nerve\textquotedblright{}.

\begin{quote}
\hi{बूढ़े} \hi{हाथों} \hi{पर} \hi{नसें} \hi{साफ़} \hi{दिखती} \hi{हैं।}
\end{quote}

\begin{quote}
\emph{būṛhe hāthoṁ par naseṁ sāf dikhtī haiṁ.}
\end{quote}

\begin{quote}
The veins show clearly on old hands.
\end{quote}

Feminine.

\begin{cousinweb}
\textbf{\hi{नस}} (\emph{nas}) is a tadbhava: CDIAL lists Hindi \emph{nas}, \textquotedblleft{}muscle, tendon, vein, nerve\textquotedblright{}, under Sanskrit \emph{snasā}, \textquotedblleft{}sinew\textquotedblright{} (CDIAL 13784). Everyday Hindi uses one word for every cord under the skin.
\end{cousinweb}

\begin{grammarlens}[title={What you've built}]
One more word for describing what you lived and felt.
\end{grammarlens}

\subsection*{Your turn}
\begin{itemize}
\raggedright
  \item \emph{Say it:} \emph{nas}
  \item \emph{Say it:} \emph{nas}, once more
  \item \emph{Say it:} \emph{būṛhe hāthoṁ par naseṁ sāf dikhtī haiṁ.} — The veins show clearly on old hands.
  \item \emph{Recall:} say the Hindi for a lung, then the Hindi for the brain, then say \emph{nas} again
\end{itemize}

\begin{tcolorbox}[breakable,colback=teal!4,colframe=teal!35!black,title={Before you move on}]
What does \hi{नस} mean? (\textquotedblleft{}A vein; a nerve\textquotedblright{}.) Say it once more.
\end{tcolorbox}
\section[\texorpdfstring{\hi{पसली} (paslī)}{पसली (paslī)}]{\hi{पसली} (paslī) — a rib}

\label{lesson:HI-C478-pasli}



\begin{quote}
Before the new one: say the Hindi for the brain, then the Hindi for a vein.
\end{quote}

\subsection*{You'll want to know: \hi{पसली}}
\textbf{\hi{पसली}} — \emph{paslī} — \textquotedblleft{}a rib\textquotedblright{}.

\begin{quote}
\hi{गिरने} \hi{से} \hi{दादा} \hi{की} \hi{एक} \hi{पसली} \hi{टूट} \hi{गई।}
\end{quote}

\begin{quote}
\emph{girne se dādā kī ek paslī ṭūṭ gaī.}
\end{quote}

\begin{quote}
Grandfather broke a rib in a fall.
\end{quote}

Feminine.

\begin{cousinweb}
\textbf{\hi{पसली}} (\emph{paslī}) is a tadbhava: CDIAL lists Hindi \emph{paslī}, \textquotedblleft{}rib\textquotedblright{}, under Sanskrit \emph{párśu}, \textquotedblleft{}rib\textquotedblright{} (CDIAL 7948).
\end{cousinweb}

\begin{grammarlens}[title={What you've built}]
One more word for describing what you lived and felt.
\end{grammarlens}

\subsection*{Your turn}
\begin{itemize}
\raggedright
  \item \emph{Say it:} \emph{paslī}
  \item \emph{Say it:} \emph{paslī}, once more
  \item \emph{Say it:} \emph{girne se dādā kī ek paslī ṭūṭ gaī.} — Grandfather broke a rib in a fall.
  \item \emph{Recall:} say the Hindi for the brain, then the Hindi for a vein, then say \emph{paslī} again
\end{itemize}

\begin{tcolorbox}[breakable,colback=teal!4,colframe=teal!35!black,title={Before you move on}]
What does \hi{पसली} mean? (\textquotedblleft{}A rib\textquotedblright{}.) Say it once more.
\end{tcolorbox}
\section[\texorpdfstring{\hi{गुर्दा} (gurdā)}{गुर्दा (gurdā)}]{\hi{गुर्दा} (gurdā) — a kidney}

\label{lesson:HI-C478-gurda}



\begin{quote}
Before the new one: say the Hindi for a vein, then the Hindi for a rib.
\end{quote}

\subsection*{You'll want to know: \hi{गुर्दा}}
\textbf{\hi{गुर्दा}} — \emph{gurdā} — \textquotedblleft{}a kidney\textquotedblright{}.

\begin{quote}
\hi{हर} \hi{आदमी} \hi{के} \hi{दो} \hi{गुर्दे} \hi{होते} \hi{हैं।}
\end{quote}

\begin{quote}
\emph{har ādmī ke do gurde hote haiṁ.}
\end{quote}

\begin{quote}
Every person has two kidneys.
\end{quote}

Masculine.

\begin{cousinweb}
\textbf{\hi{गुर्दा}} (\emph{gurdā}) is usually traced to Persian \emph{gurda}, \textquotedblleft{}kidney\textquotedblright{}.
\end{cousinweb}

\begin{grammarlens}[title={What you've built}]
One more word for describing what you lived and felt.
\end{grammarlens}

\subsection*{Your turn}
\begin{itemize}
\raggedright
  \item \emph{Say it:} \emph{gurdā}
  \item \emph{Say it:} \emph{gurdā}, once more
  \item \emph{Say it:} \emph{har ādmī ke do gurde hote haiṁ.} — Every person has two kidneys.
  \item \emph{Recall:} say the Hindi for a vein, then the Hindi for a rib, then say \emph{gurdā} again
\end{itemize}

\begin{tcolorbox}[breakable,colback=teal!4,colframe=teal!35!black,title={Before you move on}]
What does \hi{गुर्दा} mean? (\textquotedblleft{}A kidney\textquotedblright{}.) Say it once more.
\end{tcolorbox}
